% slides.tex --- minimal slide deck distilled from main.tex
% One page per slide (Beamer). Build with:  latexmk -pdf slides.tex
% Each PDF page is a self-contained slide: title + a few equations
% (or one equation with a short proof below). Insert pages into
% PowerPoint as images, or copy the rendered math directly.

\documentclass[aspectratio=169,12pt]{beamer}

\newif\ifdarkmode
\darkmodetrue  % Uncomment this line to enable dark mode

\usetheme{default}
\usecolortheme{default}
\setbeamertemplate{navigation symbols}{}
\setbeamertemplate{footline}{}          % no page numbers
\setbeamerfont{frametitle}{size=\large,series=\bfseries}
\setbeamertemplate{frametitle}[default][left]

% Dark mode toggle
\ifdarkmode
  % Dark mode colors
  \setbeamercolor{background canvas}{bg=black}
  \setbeamercolor{normal text}{fg=white}
  \setbeamercolor{frametitle}{fg=white}
  \setbeamercolor{title}{fg=white}
  \setbeamercolor{subtitle}{fg=white}
  \setbeamercolor{author}{fg=white}
  \setbeamercolor{itemize item}{fg=white}
  \setbeamercolor{itemize subitem}{fg=white}
  \setbeamercolor{enumerate item}{fg=white}
\else
  % Light mode colors (default)
  \setbeamercolor{normal text}{fg=black}
  \setbeamercolor{frametitle}{fg=black}
  \setbeamercolor{title}{fg=black}
  \setbeamercolor{subtitle}{fg=black}
  \setbeamercolor{author}{fg=black}
  \setbeamercolor{itemize item}{fg=black}
  \setbeamercolor{itemize subitem}{fg=black}
  \setbeamercolor{enumerate item}{fg=black}
\fi
\usebeamercolor[fg]{normal text}

\usepackage[utf8]{inputenc}
\usepackage[T1]{fontenc}
\usepackage{amsmath, amssymb, amsthm}
\usepackage{amsfonts}
\usepackage{bm}
\usepackage{booktabs}
\usepackage{graphicx}

% --- figure-path macros copied verbatim from main.tex ---
\newcommand{\synfig}[1]{notebooks/figures/synthetic/gene_expression_synthetic/imaging_synthetic/#1}
\newcommand{\kircfig}[1]{notebooks/figures/tcga_kirc/gene_expression/organ_radiomics/#1}
\newcommand{\nsclcfig}[1]{notebooks/figures/nsclc/gene_expression/organ_radiomics_subset/#1}

% --- macros copied verbatim from main.tex so notation matches ---
\newcommand{\R}{\mathbb{R}}
\newcommand{\E}{\mathbb{E}}
\newcommand{\N}{\mathcal{N}}
\newcommand{\Cov}{\operatorname{Cov}}
\newcommand{\Var}{\operatorname{Var}}
\newcommand{\rank}{\operatorname{rank}}
\newcommand{\tr}{\operatorname{tr}}
\newcommand{\Sg}{\Sigma_{G}}
\newcommand{\Sx}{\Sigma_{X}}
\newcommand{\Sgx}{\Sigma_{GX}}
\newcommand{\Sxg}{\Sigma_{XG}}
\newcommand{\Sgcx}{\Sigma_{G\mid X}}
\newcommand{\T}{^{\top}}

\title{Image-Identifiable Genomic Subspaces}
\subtitle{A Linear--Gaussian Model of Radiogenomic Recoverability}
\author{Joseph Rich \and Lior Pachter}
\date{}

\begin{document}

\begin{frame}
  \titlepage
\end{frame}

% ------------------------------------------------------------------
% Slide: Notation table
% ------------------------------------------------------------------
\begin{frame}{Setup and notation}
  \small
  \centering
  \renewcommand{\arraystretch}{1.25}
  \begin{tabular}{ll}
    \toprule
    $G \in \R^p,\ X \in \R^d$ & genomic / imaging feature vectors\\
    $\mathbf{G}\in\R^{n\times p},\ \mathbf{X}\in\R^{n\times d}$ & stacked data ($n$ patients)\\
    $A \in \R^{d\times p}$ & genomics-to-imaging transfer map\\
    $\sigma^2$ & isotropic imaging noise variance\\
    $\Sg,\ \Sx,\ \Sgx$ & genomic / imaging / cross covariance\\
    $\Sgcx$ & posterior covariance (recoverability floor)\\
    $B = \Sgx \Sx^{-1}$ & Bayes-optimal map $x \mapsto \E[g\mid x]$\\
    $M = \sigma^{-1} A \Sg^{1/2}$ & whitened transfer operator; $d_i$ its sing.\ values\\
    $\rho_i,\ R_i = \rho_i^2$ & $i$-th canonical correlation / recoverability\\
    $a_i,\ b_i$ & $i$-th canonical genomic / imaging direction\\
    $I(G;X)$ & mutual information between modalities\\
    $k$ & dimension of image-identifiable genomic subspace\\
    \bottomrule
  \end{tabular}
\end{frame}

% ------------------------------------------------------------------
% Slide: Generative model
% ------------------------------------------------------------------
\begin{frame}{The linear--Gaussian generative model}
  \begin{align*}
    G &\sim \N(0, \Sg),\\[2pt]
    X \mid G &\sim \N(A G,\ \sigma^2 I_d), \qquad A \in \R^{d \times p}.
  \end{align*}
  \vfill
  Equivalently,
  \[
    X = A G + \varepsilon, \qquad
    \varepsilon \sim \N(0, \sigma^2 I_d), \quad \varepsilon \perp G.
  \]
  \vfill
  \footnotesize
  Simplest model with a non-trivial answer: \emph{every} downstream quantity
  is closed-form, so it acts as an exact upper bound on what any estimator can recover.
\end{frame}

% ------------------------------------------------------------------
% Slide: Joint law
% ------------------------------------------------------------------
\begin{frame}{Joint Gaussian law}
  \[
  \begin{pmatrix} G\\ X \end{pmatrix}
  \sim
  \N\!\left(
  \mathbf{0},
  \begin{pmatrix}
  \Sg & \Sg A\T\\
  A\Sg & A\Sg A\T + \sigma^2 I_d
  \end{pmatrix}
  \right).
  \]
  \vfill
  In particular,
  \[
    X \sim \N(0,\Sx), \qquad \Sx = A\Sg A\T + \sigma^2 I_d,
  \]
  \[
    \Sgx = \Cov(G,X) = \Sg A\T.
  \]
  \vfill
  \footnotesize
  Imaging covariance splits additively into a genomics-driven part $A\Sg A\T$
  and an irreducible part $\sigma^2 I_d$.
\end{frame}

% ------------------------------------------------------------------
% Slide: Joint law -- proof
% ------------------------------------------------------------------
\begin{frame}{Joint law --- proof sketch}
  \small
  $(G,X)$ is an affine map of independent Gaussians $(G,\varepsilon)$, hence jointly Gaussian; both are centered, so the mean is $0$.
  \vfill
  \begin{align*}
    \Cov(G,G) &= \Sg,\\
    \Cov(G,X) &= \Cov(G,\,AG+\varepsilon) = \Cov(G,AG) = \Sg A\T,\\
    \Cov(X,X) &= \Cov(AG) + \Cov(\varepsilon) = A\Sg A\T + \sigma^2 I_d.
  \end{align*}
  \vfill
  Cross terms vanish by $G \perp \varepsilon$. Collecting the blocks gives the joint covariance. \hfill$\qed$
\end{frame}

% ------------------------------------------------------------------
% Slide: Posterior
% ------------------------------------------------------------------
\begin{frame}{Bayes-optimal recovery: the posterior}
  $G \mid X \sim \N(\hat g(X),\ \Sgcx)$ with
  \begin{align*}
    \hat g(X) &= \E[G\mid X] = \Sg A\T\big(A\Sg A\T + \sigma^2 I_d\big)^{-1} X = BX,\\[2pt]
    \Sgcx &= \Sg - \Sg A\T\big(A\Sg A\T + \sigma^2 I_d\big)^{-1} A\Sg.
  \end{align*}
  \vfill
  Precision form (Woodbury):
  \[
    \Sgcx^{-1} = \Sg^{-1} + \sigma^{-2} A\T A,
    \qquad B = \sigma^{-2}\,\Sgcx\, A\T.
  \]
  \vfill
  \footnotesize
  $\Sgcx$ is a \textbf{hard floor} on residual genomic uncertainty: no algorithm
  can do better. Imaging \emph{adds} precision $\sigma^{-2}A\T A$ to the prior.
\end{frame}

% ------------------------------------------------------------------
% Slide: Mutual information
% ------------------------------------------------------------------
\begin{frame}{Mutual information and the SNR decomposition}
  \[
    I(G;X)
    = \tfrac12 \log \frac{\det \Sg}{\det \Sgcx}
    = \tfrac12 \log \det\!\big(I_d + \sigma^{-2} A\Sg A\T\big).
  \]
  \vfill
  With $\lambda_1 \ge \cdots \ge \lambda_r > 0$ the nonzero eigenvalues of $A\Sg A\T$:
  \[
    I(G;X) = \tfrac12 \sum_{i=1}^{r} \log\!\Big(1 + \tfrac{\lambda_i}{\sigma^2}\Big).
  \]
  \vfill
  \footnotesize
  Each image-visible genomic direction contributes $\tfrac12\log(1+\mathrm{SNR}_i)$.
  Directions with small signal relative to noise contribute nothing.
\end{frame}

% ------------------------------------------------------------------
% FIGURE: MI per direction (synthetic)
% ------------------------------------------------------------------
\begin{frame}{A handful of directions carry almost all the information}
  \centering
  \includegraphics[width=0.72\linewidth]{\synfig{mi_per_direction.pdf}}
  \par\footnotesize
  Per-direction $\tfrac12\log(1+d_i^2)$ on a synthetic instance ($n{=}800$, $p{=}80$, $d{=}40$);
  cumulative MI saturates at $\approx 6.45$ bits after the first few components.
\end{frame}

% ------------------------------------------------------------------
% Slide: Recoverability = squared canonical correlation
% ------------------------------------------------------------------
\begin{frame}{Recoverability $=$ squared canonical correlation}
  Whiten $U = \Sg^{-1/2}G$, transfer operator $M = \sigma^{-1}A\Sg^{1/2} = \sum_i d_i u_i v_i\T$.
  Canonical direction $a_i = \Sg^{-1/2} v_i$, score $S_i = a_i\T G$:
  \[
    \Var(S_i \mid X) = \frac{1}{1+d_i^2},
    \qquad
    \rho_i = \frac{d_i}{\sqrt{1+d_i^2}}.
  \]
  \vfill
  \[
    \boxed{\;R(a_i) = 1 - \Var(S_i\mid X) = \frac{d_i^2}{1+d_i^2} = \rho_i^2.\;}
  \]
  \vfill
  \footnotesize
  The whole analysis reduces to CCA between $\mathbf{G}$ and $\mathbf{X}$:
  only the moments $\Sg,\Sx,\Sgx$ enter --- $A$ and $\sigma^2$ never appear individually.
\end{frame}

% ------------------------------------------------------------------
% Slide: Named direction recoverability score
% ------------------------------------------------------------------
\begin{frame}{Recoverability of a named genomic direction}
  For a fixed score $Y = v\T G$ (a gene, pathway, polygenic score):
  \[
    R(v) = 1 - \frac{v\T \Sgcx\, v}{v\T \Sg\, v} \;\in\; [0,1].
  \]
  \vfill
  Equals the population $R^2$ of the Bayes-optimal predictor of $Y$ from $X$. The core matrix has a moment-only form:
  \[
    \Sg - \Sgcx = \Sg A\T \Sx^{-1} A \Sg = \Sgx\,\Sx^{-1}\,\Sxg.
  \]
  \vfill
  \footnotesize
  $R(v)=0$: imaging carries no information about $Y$. \quad
  $R(v)=1$: imaging determines it exactly.
\end{frame}

% ------------------------------------------------------------------
% Slide: Eigenproblem = CCA
% ------------------------------------------------------------------
\begin{frame}{The recoverability eigenproblem $=$ CCA}
  \[
    \max_{v \neq 0}\; R(v)
    = \max_{v \neq 0}\;
    \frac{v\T \big(\Sgx \Sx^{-1} \Sxg\big) v}{v\T \Sg\, v}
  \]
  has stationary points solving the generalized eigenproblem
  \[
    \Sgx\, \Sx^{-1}\, \Sxg\; v = \mu\, \Sg\, v,
    \qquad \mu_i = \rho_i^2,\quad v = a_i.
  \]
  \vfill
  \footnotesize
  Solutions are exactly the canonical correlation directions, and
  $\max_v R(v) = \rho_1^2$. Deliverable: a \emph{recoverability spectrum}
  $\rho_1^2 \ge \rho_2^2 \ge \cdots$ with genomic directions $\{a_i\}$.
\end{frame}

% ------------------------------------------------------------------
% Slide: Rank limit
% ------------------------------------------------------------------
\begin{frame}{Rank limit and identifiability}
  \[
    \rank(\Sg - \Sgcx) = \rank(A\Sg A\T)
    \le \min\!\big(d,\ p,\ \rank A,\ \rank \Sg\big).
  \]
  \vfill
  \footnotesize
  \begin{itemize}
    \item At most $\min(d,p,\rank A,\rank\Sg)$ canonical correlations are nonzero.
    \item Even with $p=20{,}000$ genes, imaging exposes only a projection whose
      dimension is capped by the channel --- not the gene panel size.
    \item Imaging features $x = \phi(\text{image})$ are themselves lossy:
      post-processing can only shrink $\rank(\Sgx)$ and lower each $\rho_i$
      (data-processing inequality).
  \end{itemize}
\end{frame}

% ------------------------------------------------------------------
% FIGURE: synthetic spectrum + truth validation
% ------------------------------------------------------------------
\begin{frame}{The estimator recovers the planted canonical structure}
  \centering
  \includegraphics[width=0.49\linewidth]{\synfig{recoverability_spectrum.pdf}}
  \hfill
  \includegraphics[width=0.49\linewidth]{\synfig{synthetic_truth_validation.pdf}}
  \par\footnotesize
  Synthetic instance with known truth: leading $\hat\rho_i^2$ matches the planted
  values (left); estimated quantities track the closed-form predictions (right).
\end{frame}

% ------------------------------------------------------------------
% Slide: Estimation regimes
% ------------------------------------------------------------------
\begin{frame}{Finite-sample regimes ($\kappa_G = p^\star/n,\ \kappa_X = d^\star/n$)}
  \footnotesize
  \begin{enumerate}
    \item \textbf{Classical} ($\kappa_G+\kappa_X \ll 1$): sample CCA consistent,
      $\hat\rho_i \to \rho_i$ at rate $n^{-1/2}$.
    \item \textbf{High-dimensional} ($\kappa_G+\kappa_X \lesssim 1$): null squared
      correlations follow a Wachter law; leading $\hat\rho_i$ inflated. Use
      regularized CCA.
    \item \textbf{Saturating} ($p^\star + d^\star \ge n$): $\hat\rho_1=\cdots=1$
      regardless of association; spectrum not estimable.
  \end{enumerate}
  \vfill
  Working rule: reduce \emph{both} modalities so $p^\star + d^\star \le n/\tau$, $\tau \in [5,10]$.
  Report cross-validated $\hat R_i^{\mathrm{cv}}$ and a scale-matched permutation null.
\end{frame}

% ------------------------------------------------------------------
% FIGURE: synthetic CV + permutation
% ------------------------------------------------------------------
\begin{frame}{In-sample recoverability is inflated; CV + permutation are honest}
  \centering
  \includegraphics[width=0.49\linewidth]{\synfig{cv_recoverability.pdf}}
  \hfill
  \includegraphics[width=0.49\linewidth]{\synfig{permutation_null.pdf}}
  \par\footnotesize
  In-sample $\hat R_i$ overshoots; $5$-fold $\hat R_i^{\mathrm{cv}}$ tracks the
  planted truth (left). Permutation null recovers the planted rank, $p<0.01$ for
  the top three (right).
\end{frame}

% ------------------------------------------------------------------
% Slide: Estimation -- regularized CCA
% ------------------------------------------------------------------
\begin{frame}{Estimation: regularized CCA}
  Shrinkage covariances (Ledoit--Wolf) $\hat\Sg, \hat\Sx$ and
  $\hat\Sgx = n^{-1}\mathbf{G}\T\mathbf{X}$. Then
  \[
    \hat M = \hat\Sg^{-1/2}\,\hat\Sgx\,\hat\Sx^{-1/2}
    = \hat U \hat D \hat V\T,
  \]
  \[
    \hat\rho_i = \hat d_i,\quad
    \hat R_i = \hat\rho_i^2,\quad
    \hat a_i = \hat\Sg^{-1/2}\hat u_i,\quad
    \hat b_i = \hat\Sx^{-1/2}\hat v_i.
  \]
  \vfill
  \footnotesize
  Gaussian-copula marginal transform $\tilde M_{ij}=\Phi^{-1}(\mathrm{rank}_i/(n+1))$
  applied to both sides --- radiomic marginals have excess kurtosis in the tens.
\end{frame}

% ------------------------------------------------------------------
% Slide: Bernoulli copula extension
% ------------------------------------------------------------------
\begin{frame}{Bernoulli extension for mutations (Gaussian copula)}
  Latent $Z \sim \N(0,\Sigma_Z)$, $\mathrm{diag}(\Sigma_Z)=I_p$;
  thresholding $G_j = \mathbf{1}\{Z_j > \tau_j\}$, $\tau_j = \Phi^{-1}(1-p_j)$.
  Imaging couples to the latent: $X \mid Z \sim \N(BZ, \sigma^2 I_d)$.
  \vfill
  Data-processing bound, tight up to the binarization gap:
  \[
    I(G;X) \le I(Z;X) = -\tfrac12 \sum_{i=1}^{k} \log\big(1-\rho_i^2\big).
  \]
  \vfill
  \footnotesize
  Per-gene AUC ceiling:
  $\mathrm{AUC}_j \le \Phi\!\big(\sqrt{R_j^Z/(1-R_j^Z)}\big/\sqrt2\big)$.
  All CCA / MI machinery transfers with $\Sg \to \Sigma_Z$.
\end{frame}

% ------------------------------------------------------------------
% Slide: KIRC results
% ------------------------------------------------------------------
\begin{frame}{Application: TCGA-KIRC ($n=190$)}
  \footnotesize
  Organ radiomics ($d_{\mathrm{native}}=107$), $2000$ HVGs, copula marginals,
  $p^\star=d^\star=38$.
  \vfill
  \begin{itemize}
    \item In-sample $\hat\rho_1 \approx 0.76$ ($\hat R_1 \approx 0.58$), $3.9$ bits.
    \item Cross-validated $\hat R_1^{\mathrm{cv}} \approx 0.07$--$0.1$; variance share $0.27 \to 0.011$.
    \item Permutation $p<0.05$ for leading-3 in $12/12$ preprocessing variants.
    \item $I(G;X)\big|_{\text{top-}3} \le 1.58$ bits (one-sided $95\%$ UCL).
    \item Effective rank $\hat k_{\mathrm{eff}} = 1$ (median).
  \end{itemize}
  \vfill
  \emph{One genomic direction is robustly image-identifiable; recoverability $\approx 0.1$.}
\end{frame}

% ------------------------------------------------------------------
% FIGURE: KIRC spectrum + MI upper bound
% ------------------------------------------------------------------
\begin{frame}{TCGA-KIRC: one weak but real image-identifiable direction}
  \centering
  \includegraphics[width=0.49\linewidth]{\kircfig{cv_recoverability.pdf}}
  \hfill
  \includegraphics[width=0.49\linewidth]{\kircfig{mi_upper_bound.pdf}}
  \par\footnotesize
  Cross-validated $\hat R_i$ is a small positive residual (left); observed
  cumulative MI sits at the upper edge of the null, $I(G;X)\le 1.58$ bits (right).
\end{frame}

% ------------------------------------------------------------------
% FIGURE: KIRC per-gene AUC
% ------------------------------------------------------------------
\begin{frame}{KIRC: recoverability score predicts which genes imaging can read}
  \centering
  \includegraphics[width=0.7\linewidth]{\kircfig{per_gene_auc.pdf}}
  \par\footnotesize
  High-recoverability genes clear the chance line (AUC $\approx 0.65$);
  low-recoverability genes sit at $0.5$. Random forest does not exceed the linear ceiling.
\end{frame}

% ------------------------------------------------------------------
% Slide: NSCLC results
% ------------------------------------------------------------------
\begin{frame}{Application: NSCLC ($n=90$) --- the counterpoint}
  \footnotesize
  Same pipeline, lung organ radiomics, $p^\star=d^\star=18$.
  \vfill
  \begin{itemize}
    \item In-sample $\hat\rho_1 = 0.82$, $2.6$ bits; null well-separated ($p=0.005$).
    \item Held-out canonical $R^2 = 0.40$; $\hat R_1^{\mathrm{cv}} = 0.33$.
    \item Per-gene AUC: identifiable $0.73$ / $0.69$ vs non-identifiable $0.51$ / $0.53$.
    \item Survives de-confounding (scanner / kernel): $\hat R_1^{\mathrm{cv}} \approx 0.24$--$0.30$.
  \end{itemize}
  \vfill
  Two cohorts, same model: borderline single direction ($\approx 0.1$) in RCC,
  solid single direction ($\approx 0.33$) in NSCLC --- same low rank, different strength.
\end{frame}

% ------------------------------------------------------------------
% FIGURE: NSCLC spectrum + per-gene AUC
% ------------------------------------------------------------------
\begin{frame}{NSCLC: a solid single direction, an order of magnitude above KIRC}
  \centering
  \includegraphics[width=0.49\linewidth]{\nsclcfig{cv_recoverability.pdf}}
  \hfill
  \includegraphics[width=0.49\linewidth]{\nsclcfig{per_gene_auc.pdf}}
  \par\footnotesize
  Leading $\hat R_1^{\mathrm{cv}}\approx 0.33$ (left); per-gene AUC $0.73/0.69$ for
  identifiable vs $0.51/0.53$ for non-identifiable genes (right).
\end{frame}

% ------------------------------------------------------------------
% FIGURE: KIRC vs NSCLC sample-complexity sweeps
% ------------------------------------------------------------------
\begin{frame}{Honest signal lives below the in-sample bulk at all cohort sizes}
  \centering
  \includegraphics[width=0.49\linewidth]{\kircfig{sample_complexity_sweep.pdf}}
  \hfill
  \includegraphics[width=0.49\linewidth]{\nsclcfig{sample_complexity_sweep.pdf}}
  \par\footnotesize
  Sample-complexity sweeps, KIRC (left) and NSCLC (right): $\hat R_1^{\mathrm{cv}}$
  is visible only through CV/permutation, never the raw in-sample spectrum.
\end{frame}

% ------------------------------------------------------------------
% Slide: Takeaway
% ------------------------------------------------------------------
\begin{frame}{Takeaway}
  \begin{itemize}
    \item Radiogenomics should target \textbf{image-identifiable low-dimensional
      genomic subspaces}, not individual genes.
    \item Recoverability of canonical direction $i$ is exactly $\rho_i^2$ ---
      computable by regularized CCA from $\Sg, \Sx, \Sgx$ alone.
    \item The recoverable dimension is a property of the \emph{biology--imaging
      channel}, bounded by $\rank(A)$, not by panel size.
    \item Honest estimands are cross-validated and permutation-calibrated;
      in-sample spectra are badly inflated when $p,d \sim n$.
  \end{itemize}
\end{frame}

\end{document}
